% TOP_PROBLEM: {"id": 2728, "title": "Kirby Problem 1.69", "category_id": 7, "category": {"id": 7, "name": "topology", "display_name": "Topology", "description": "Properties preserved under continuous deformations.", "slug": "topology", "order_index": 7, "created_at": "2026-07-31T15:26:25.671Z"}, "status": "open"}
% FIRST_POSTED: null
% ENTRY_KIND: statement_only
% Imported from: batch11.tex; UnsolvedMath ID 2728
\documentclass[11pt]{article}
\usepackage[margin=25mm]{geometry}
\usepackage{fontspec}
\setmainfont{Latin Modern Roman}
\setsansfont{Latin Modern Sans}
\usepackage{amsmath,amssymb,mathtools}
\usepackage{unicode-math}
\setmathfont{Latin Modern Math}
\usepackage{xurl,hyperref}
\setcounter{secnumdepth}{0}
\setlength{\parindent}{0pt}
\setlength{\parskip}{5pt}
\setlength{\emergencystretch}{3em}
\newcommand{\sref}[2]{\hyperlink{source:#1:#2}{[S#2]}}
\newcommand{\cataloguescope}[1]{\paragraph{Recorded target.}#1}
\begin{document}
% UnsolvedMath ID 2728; source catalogue code KP-1.69
\setcounter{section}{2727}
\section{Kirby Problem 1.69}\label{2728}
{\small\sffamily UnsolvedMath ID 2728 · Topology}
\subsection{Definitions and mathematical statement}

\paragraph{Shared notation.}$\mathbb N=\{1,2,\ldots\}$, and $\mathbb Z,\mathbb Q,\mathbb R,\mathbb C$ denote the integers, rationals, reals and complex numbers. Any other conventions are specified locally.


An oriented link $L\subset(S^3,\xi_{\mathrm{std}})$ is positively transverse when its oriented tangent is positively transverse to the standard contact planes. For a Seifert surface $\Sigma\subset S^3$, the self-linking $\operatorname{sl}_\Sigma(L)$ is the linking with a contact push-off defined by a nonzero contact-plane section over $\Sigma$. A Seifert surface is a compact oriented embedded spanning surface without closed components. A quasipositive braid is a product of conjugates of positive braid generators. A strongly quasipositive braid is a product of positive embedded band generators
\[\sigma_{i,j}=(\sigma_i\cdots\sigma_{j-2})\sigma_{j-1}(\sigma_i\cdots\sigma_{j-2})^{-1},\qquad i<j.\]
Their closures define quasipositive and strongly quasipositive link types respectively. A quasipositive Seifert surface is the corresponding disk-and-positive-band surface for a strongly quasipositive braid, up to isotopy. In a general contact 3-manifold $(Y,\xi)$ the relative self-linking is defined with the same spanning-surface section convention. A Legendrian ribbon is an oriented surface retracting onto a Legendrian graph, tangent to $\xi$ along the graph and transverse to $\xi$ elsewhere, with positively transverse boundary.
\paragraph{Statement recorded in the supplied source.}
If a transverse link $L\subset(S^3,\xi_{\mathrm{std}})$ admits a Seifert surface $\Sigma$ with $\operatorname{sl}_\Sigma(L)=-\chi(\Sigma)$, must $L$ be strongly quasipositive? Must this particular surface $\Sigma$ be isotopic to a quasipositive Seifert surface? More generally, for a transverse link in a contact 3-manifold $(Y,\xi)$ satisfying $\operatorname{sl}_{\xi,\Sigma}(L)=-\chi(\Sigma)$, must $\Sigma$ be isotopic to a Legendrian ribbon? \sref{2728}{2}

\subsection{Short English statement}
Does sharpness of the Seifert-surface Bennequin bound imply strong quasipositivity, including quasipositivity of the surface itself? Ask the analogous Legendrian-ribbon question in general contact 3-manifolds.
\subsection{Sources}
\begin{description}
\item[\hypertarget{source:2728:1}{S1}] ULAM.AI and the UnsolvedMath Contributors, UnsolvedMath: A Curated Collection of Open Mathematics Problems, record 2728 (KP-1.69). CC BY 4.0. \url{https://huggingface.co/datasets/ulamai/UnsolvedMath}.
\item[\hypertarget{source:2728:2}{S2}] R. İnanç Baykur, Robion C. Kirby and Daniel Ruberman, eds., \emph{K3: A New Problem List in Low-Dimensional Topology}, Mathematical Surveys and Monographs 295, American Mathematical Society (2026), Problem 1.69 and its remarks; Chapters 1 and 2 for the stated conventions. \url{https://math.berkeley.edu/sites/default/files/surv-295-ruberman-watermarked-author-pdf.pdf}.
\end{description}
\label{end:2728}
\end{document}
